\subsubsection{0D Outlet and Network Models}
\label{subsubsec:0d-models}

A 0D model is a lumped ODE or circuit model for pressures, flows, volumes, and
loads. It retains compartment pressures, branch flows, resistances,
compliances, and, when the circuit includes inertial effects, inertances. The
model preserves a time-dependent pressure--flow relation for parts of the
circulation that are not spatially resolved. Windkessel and RCR models are the
canonical terminal examples: they supply an outlet impedance or downstream load
for a local 1D or 3D vessel solve
\cite{ShiEtAl2011BloodFlowModelReview,KopplHelmig2023DimensionReducedModeling}.

A 0D model closes the circulation seen by the local vessel model; it does not
resolve the local stenosis. The spatial coordinate has been integrated out, so
the model has no local stenosis geometry, no axial pressure-drop field, no
throat acceleration profile, and no recirculation region. It also has no
resolved wall-shear stress or wall traction, because it does not carry a lumen
surface, a stress tensor on that surface, or a wall-normal output operator.

The 0D tier is therefore useful boundary and network context for this report. It
can prescribe an outlet load, organize distal-bed assumptions, or provide a
lumped circulation closure. It is not the selected local stenosis-resolution
tier and cannot replace the distributed area--flow state used in the numerical
study in this report.
